\infographicchapter{Myelography}{myelography.png}
\label{ch:infographic-myelography}

\section{Overview}

Myelography is a well-established medical test that looks closely at the spinal canal, spinal cord, and the nerves branching off the spine. Unlike a routine X-ray or MRI, it involves a minor spinal puncture. After numbing the skin and cleaning the area carefully, a radiologist uses live X-ray guidance (fluoroscopy) to place a very thin needle into the fluid-filled space around the spinal cord---usually in the lower back, or occasionally the neck. A special water-based contrast dye is injected into the spinal fluid (cerebrospinal fluid, or CSF). Gently tilting the exam table lets gravity carry the dye along the spine so live X-rays and a detailed CT scan (called a CT myelogram) can clearly outline the spinal cord and nerves.

Although MRI is usually the first test doctors order for spine problems, myelography remains vital when an MRI cannot be done---such as for patients with certain older pacemakers, metal fragments, or severe claustrophobia---or when metal spine screws and rods blur an MRI. Common reasons for myelography include spinal canal narrowing (stenosis), slipped discs, pinched nerves, checking a spine after hardware surgery, evaluating spinal tumors or nerve inflammation, and finding the exact spot of a spinal fluid leak.

Key benefits include the ability to take pictures while your spine moves (such as bending forward or backward), crisp detail where bone meets nerves, and much less blurring from surgical metal than MRI. On the other hand, it is an invasive procedure that uses X-ray radiation and cannot see inside the spinal cord tissue itself as well as MRI. The most common side effect is a temporary headache after the spinal puncture, while rarer risks include bleeding, infection, allergic dye reactions, or seizures. Your care team will carefully check your blood thinners, allergies, and pregnancy status beforehand.

\section{Underlying Physical Principles}
\begin{itemize}
    \item \textbf{Highlighting the Spinal Fluid With Contrast Dye:} Normal spinal fluid looks clear on an X-ray. Adding a safe, water-based iodine dye makes the fluid absorb X-rays strongly, turning the fluid sac bright white so the darker spinal cord and branching nerves stand out in sharp relief.
    \item \textbf{Guiding the Dye With Gravity:} Because the contrast dye is slightly heavier than natural spinal fluid, gently tilting the padded exam table uses gravity to glide the dye smoothly into the lower back, middle back, or neck regions as needed.
    \item \textbf{Live X-Ray Guidance (Fluoroscopy):} Low-dose, pulsed live X-ray video lets the radiologist guide the thin spinal needle safely into place, watch the dye flow in real time, and capture pictures while the spine is in different positions.
    \item \textbf{High-Detail CT Cross-Sections:} Right after the X-ray portion, a fast CT scanner takes thin 3D slices of the dye-filled spine. Special computer programs sharpen the images and reduce glare around metal spine screws so doctors can clearly see bone spurs, discs, and nerve roots.
\end{itemize}

\section{Imaging Mechanism}
\begin{enumerate}
    \item \textbf{Sterile Setup and Positioning:} You lie face-down or on your side on a tilting X-ray table. The radiologist cleans the skin on your lower back (or occasionally your neck) with antiseptic, places sterile drapes, and numbs the skin with local anesthetic.
    \item \textbf{Needle Placement and Dye Injection:} Watching on a live X-ray screen, the radiologist gently guides a thin spinal needle into the fluid space around the spinal cord and injects the contrast dye. Once the needle is removed, the table tilts gently so gravity spreads the dye along the spine.
    \item \textbf{Live X-Rays and Transfer to CT:} The radiologist takes live X-ray snapshots to watch how the dye flows and outlines the nerves. You then move promptly to a nearby CT scanner, which takes detailed 3D X-ray slices while the dye is still bright.
    \item \textbf{3D Image Building and Metal Cleanup:} The CT computer builds thin, sharp slices of your spine from multiple angles and uses special software to minimize streaks around any surgical metal hardware.
    \item \textbf{Expert Review:} A radiologist studies both the live X-ray snapshots and the 3D CT myelogram slices to check the spinal fluid sac, the outline of the spinal cord, and every exiting nerve root.
\end{enumerate}

\section{Key Technical Parameters}
\noindent
\begin{tabularx}{\textwidth}{@{}l X@{}}
\toprule
\textbf{Parameter} & \textbf{Typical Value / Description} \\
\midrule
Physical Principle & Live X-ray-guided injection of contrast dye into spinal fluid to outline the spinal cord and nerves, followed by CT \\
Energy / Signal Source & Pulsed live X-ray fluoroscopy and a helical CT scanner \\
Typical Frequency / Energy & 70--110 kVp (live X-ray); 100--140 kVp (CT myelogram) \\
Spatial Resolution & Very sharp bone and nerve-outline detail (0.3--0.6 mm CT slices; 0.2--0.5 mm X-ray views) \\
Typical Scan Duration & 45--60 minutes total (30--45 minutes for puncture and X-rays, plus 10--15 minutes for CT) \\
Contrast Agent & Water-based, non-ionic iodine contrast dye approved specifically for spinal fluid injection (such as iohexol or iopamidol) \\
Ionizing Radiation & Yes; typically 5--15 mSv combined (live X-ray plus targeted spine CT) \\
\bottomrule
\end{tabularx}

\section{Clinical Usages}

\subsection{Primary Usages}
\begin{itemize}
    \item \textbf{When MRI Cannot Be Done:} Serving as the main alternative to MRI for checking spinal canal narrowing, slipped discs, and pinched nerves in people with MRI-unsafe pacemakers, metal fragments, or severe claustrophobia.
    \item \textbf{Checking Spines With Metal Hardware:} Looking for ongoing or returning nerve pain after spine surgery when metal screws, rods, or cages create too much distortion on an MRI.
    \item \textbf{Finding Spinal Fluid (CSF) Leaks:} Pinpointing the exact tear or abnormal vein connection where spinal fluid is leaking in patients who have severe headaches when standing up (spontaneous intracranial hypotension).
    \item \textbf{Pinched or Torn Nerve Roots:} Showing subtle nerve root pinching or torn nerve sleeves after neck or shoulder injuries when an MRI does not give a clear answer.
    \item \textbf{Spinal Canal Masses and Nerve Clumping:} Outlining tumors inside the spinal fluid sac, fluid cysts, clumped nerves from inflammation (arachnoiditis), or enlarged abnormal blood vessels along the spine.
\end{itemize}

\subsection{Secondary Usages}
\begin{itemize}
    \item \textbf{Movement and Weight-Bearing Views:} Taking pictures while you bend forward, lean back, or stand upright to uncover spine narrowing or slipping that hides when you lie flat for an MRI.
    \item \textbf{Planning Complex Spine Surgery:} Giving surgeons crystal-clear views of hard bone spurs versus soft disc material before delicate spine decompression or revision surgery.
    \item \textbf{Testing Spinal Fluid at the Same Time:} Collecting a small sample of spinal fluid through the same needle to test for infection, inflammation, or abnormal cells.
    \item \textbf{Solving Unclear MRI Results:} Clearing up confusion when your nerve symptoms do not match what a standard MRI shows.
\end{itemize}

\section{Advantages}
\begin{itemize}
    \item \textbf{Reliable Alternative to MRI:} Delivers sharp pictures of the spinal canal and nerves when an MRI is unsafe or impossible due to implanted devices or severe claustrophobia.
    \item \textbf{Clear Views Around Surgical Metal:} When paired with CT, shows the spine and nerves around metal screws, rods, and cages far more clearly than MRI.
    \item \textbf{Sharp Bone and CSF Leak Detail:} Clearly distinguishes hard bone spurs from soft discs and is the gold standard for pinpointing fast spinal fluid leaks.
    \item \textbf{Captures Spine Movement:} Lets doctors see how the spinal canal narrows when you bend, stretch, or bear weight---something a standard flat MRI cannot show.
    \item \textbf{Two Tests in One Step:} Allows doctors to collect a spinal fluid sample for lab testing during the same needle placement.
\end{itemize}

\section{Disadvantages}
\begin{itemize}
    \item \textbf{Risks of a Spinal Puncture:} Because a needle enters the spinal fluid space, it carries risks of a post-puncture headache (the most common side effect), bleeding around the spine, infection, nerve irritation, or a small temporary fluid leak.
    \item \textbf{Cannot See Inside the Spinal Cord:} The dye outlines the outside edge of the spinal cord and nerves, but cannot show problems inside the cord tissue itself (such as swelling, multiple sclerosis plaques, or internal cord tumors) as well as MRI.
    \item \textbf{Ionizing Radiation Exposure:} Uses X-rays for both the live guidance and the CT scan, so it is generally avoided during pregnancy.
    \item \textbf{Rare Dye Reactions:} Spinal contrast dye carries rare risks of allergic reactions, nerve lining irritation, brief confusion, or seizures. Care teams carefully check your allergies and medications beforehand.
    \item \textbf{Not Safe With Certain Bleeding or Pressure Risks:} Cannot be performed safely if you have a severe bleeding disorder, cannot pause high-risk blood thinners, have a skin infection over the back, or have dangerous pressure buildup inside the skull.
\end{itemize}

\section{Contraindications}

\subsection{Absolute}
\begin{itemize}
    \item \textbf{Skin Infection Over the Puncture Site:} Any active skin or tissue infection where the needle needs to go, because it could carry germs into the spinal fluid and cause meningitis or an abscess.
    \item \textbf{High Pressure Inside the Skull:} Suspected brain swelling or a brain mass causing high pressure, where removing spinal fluid could cause dangerous shifting of brain tissue.
    \item \textbf{Severe Uncontrolled Bleeding Risk:} Severe bleeding disorders, very low platelets, or an inability to safely pause strong blood thinners, which could lead to bleeding inside the spinal canal.
\end{itemize}

\subsection{Relative}
\begin{itemize}
    \item \textbf{Severe Dye Allergy or Seizure-Risk Medications:} A history of severe allergic reactions to iodine dye (which may require steroid premedication or a different test), or taking medications that raise seizure risk unless safely adjusted beforehand.
    \item \textbf{Pregnancy and Complex Spine Anatomy:} Usually avoided during pregnancy because of radiation to the spine and pelvis. Severe spine curvature or prior scarring may require specialized image-guided planning.
\end{itemize}

\section{Patient Safety Guidelines}
\begin{itemize}
    \item \textbf{Checking Bleeding and Brain Pressure Risks:} Your care team reviews blood thinners, platelet counts, and clotting labs before the test to prevent spinal bleeding, and checks for any signs of high pressure inside the head before doing a spinal puncture.
    \item \textbf{Using Approved Spinal Dye and Reviewing Medications:} Only non-ionic, water-based iodine dyes specifically approved for spinal fluid use are given, within strict dose limits. Staff also screen for dye allergies and check medications that can lower your seizure threshold (such as certain antidepressants or antipsychotics) to see if a temporary pause is needed.
    \item \textbf{Strict Cleanliness and Live Guidance:} The puncture is done using sterile gloves, drapes, and skin antiseptic under live X-ray guidance, always avoiding infected skin to prevent infection.
    \item \textbf{Keeping Radiation Low:} Radiologists use short, pulsed X-rays and tight beams to keep radiation as low as possible, and check for pregnancy in patients who could be pregnant.
    \item \textbf{Recovery, Fluids, and Head Elevation:} Afterward, staff watch you for headache, nausea, new weakness, or rare seizures. Drinking plenty of fluids and keeping your head propped up at a 30--45 degree angle for several hours helps keep the dye from drifting up toward the head and lowers headache risk.
\end{itemize}

\section{Prerequisite Tests}
\begin{itemize}
    \item \textbf{Blood-Clotting and Medication Review:} Because a spinal puncture carries a small risk of bleeding, your doctor checks blood thinners and may order clotting blood tests (platelets, INR/PT, aPTT) to confirm it is safe to proceed.
    \item \textbf{Neurological Exam and Prior Scans:} A brief nerve exam and review of past brain or spine scans help make sure there is no dangerous pressure buildup in the head and guide the exact spinal levels to focus on.
    \item \textbf{Allergy and Medication Screening:} Staff check for any past reactions to iodine contrast dye and review medications that can increase seizure risk so premedication or temporary medication adjustments can be arranged safely.
    \item \textbf{Pregnancy and Skin Check:} Patients who could be pregnant are screened (with a pregnancy test if needed), and the doctor inspects the skin on your back to make sure there is no rash or infection at the needle site.
\end{itemize}

\section{Preparing for the Scan}
\begin{itemize}
    \item \textbf{Drink Extra Fluids and Follow Eating Rules:} Start drinking extra water the day before your test to help prevent a post-puncture headache. Follow your center's food instructions---usually no solid food for 4--6 hours before the procedure, while clear liquids are often allowed unless sedation is planned.
    \item \textbf{Pause Certain Medications Only With Doctor Approval:} Follow your doctor's exact schedule for pausing blood thinners or seizure-threshold medications (such as certain antidepressants or antipsychotics). Never stop blood thinners or mental-health medications on your own without your prescribing doctor's guidance.
    \item \textbf{Allergy Premedication and Consent:} If you have a past iodine dye allergy, take any prescribed steroid and antihistamine pills beforehand. You will also review the steps, risks, and alternatives with your doctor and sign a consent form.
    \item \textbf{Clean Skin and Simple Clothing:} Shower with mild soap before arriving, wear loose comfortable clothes, and remove jewelry, belts, and metal items before changing into a hospital gown.
    \item \textbf{Arrange a Ride Home:} Plan to have a responsible adult drive you home after your 2--4 hour recovery stay. You will need to avoid driving and heavy activity for at least 24 hours.
\end{itemize}

\section{During Scan}
\begin{itemize}
    \item \textbf{Positioning and Table Tilting:} You will lie face-down or on your side on a padded table with cushions for support while the puncture is done. After the dye is placed, padded shoulder and foot supports keep you secure while the table gently tilts head-down or upright so gravity moves the dye. You then lie on your back for the CT scan.
    \item \textbf{Breathing and Staying Still:} Breathe slowly and steadily during the needle placement, and try not to cough, sneeze, or shift suddenly. During the X-ray snapshots and CT scan, follow simple prompts to hold your breath briefly.
    \item \textbf{What You Will Feel:} You will feel cool cleansing liquid on your back, a brief pinch and sting from the numbing medicine, and deep pressure---or an occasional brief tingle down a leg---as the thin needle enters the fluid space. The dye itself may feel slightly cool or cause mild fullness in your lower back.
    \item \textbf{When to Speak Up:} Tell the radiologist right away if you feel sharp shooting pain down an arm or leg, a sudden severe headache, nausea, dizziness, itching, or trouble breathing at any point.
    \item \textbf{How Long It Takes:} The puncture and live X-rays usually take 30--45 minutes, followed by a 10--15 minute CT scan---about 45--60 minutes in total before moving to the recovery area.
\end{itemize}

\section{Normal Appearance}
\begin{itemize}
    \item \textbf{Smooth, Even Dye Column:} The bright white contrast dye fills the spinal fluid sac smoothly and evenly without pinching, blockages, or leaks outside the sac.
    \item \textbf{Centered Spinal Cord and Free Nerves:} The spinal cord shows up as a smooth, centered darker outline ending near the upper lower back (around L1--L2), with loose, even bundles of lower nerves floating freely below it.
    \item \textbf{Even Nerve Sleeves:} The tiny dye-filled pockets around each exiting nerve root fill smoothly and evenly on both sides without being cut off or pushed aside.
\end{itemize}

\section{Observing for Abnormal Patterns}
\begin{itemize}
    \item \textbf{Pinching From Outside the Fluid Sac:} A smooth indentation or hourglass squeeze pushing into the bright dye column points to pressure from outside the sac---such as a slipped disc, bone spurs, thickened spinal ligaments, bleeding, or a bone tumor.
    \item \textbf{Masses Inside the Fluid Sac:} A sharp, rounded gap inside the dye column that pushes the spinal cord to the side---often forming a curved cap of dye around the edge---is a classic sign of a tumor growing on a nerve or the spinal lining (such as a meningioma or schwannoma).
    \item \textbf{Swelling Inside the Spinal Cord:} Smooth widening of the spinal cord itself that thins out the dye all the way around it points to swelling, a fluid cavity (syrinx), or a tumor inside the cord.
    \item \textbf{Pinched Nerve Sleeves, Fluid Leaks, and Nerve Clumping:} A nerve pocket that stops short or fails to fill shows a pinched nerve. Dye leaking outside the sac or lighting up a nearby vein pinpoints a spinal fluid leak, while nerves stuck together along the wall point to inflammation and scarring (arachnoiditis).
    \item \textbf{Technical Look-Alikes:} If dye accidentally enters the wrong thin layer around the sac, it can look like stiff, straight lines that do not move when the table tilts. Trapped air bubbles (tiny round spots), uneven dye mixing, or glare from metal spine hardware can also mimic disease if not recognized.
\end{itemize}

\section{Dealing with Uncertainties}
\begin{itemize}
    \item \textbf{Dye Outside the Main Fluid Space:} Occasionally, contrast dye collects in a tight layer just outside the spinal fluid space, failing to flow smoothly with gravity and making the test harder to read.
    \item \textbf{Blocked Flow and Metal Glare:} Very tight spinal narrowing can block dye from flowing past a certain point. Delays before the CT scan, tiny air bubbles, or heavy glare from metal spine rods can also hide subtle nerve pinching or small fluid leaks.
    \item \textbf{Limits Inside the Cord:} Because myelography only outlines the outside of the spinal cord, problems deep inside the cord tissue cannot be seen clearly, and different causes of outside pressure (like disc fragments, bone spurs, or scar tissue) can look similar.
    \item \textbf{Adjusting Position and Special CT Modes:} If dye pools unevenly, the team can roll you onto your side, tilt the table gently, take bending views, or use special metal-reducing CT software and rapid dynamic CT scans to catch fast fluid leaks.
    \item \textbf{Matching With Symptoms and Other Tests:} Doctors always match unclear pictures with your exact pain or numbness pattern and any spinal fluid lab tests. When safe, a contrast MRI---or a second puncture from another level if flow was completely blocked---may be used before surgery.
\end{itemize}

\section{Augmenting Methods}
\begin{itemize}
    \item \textbf{Spine MRI:} When safe to perform, MRI pairs with CT myelography by showing the inside of the spinal cord, bone marrow, and soft tissues in rich detail.
    \item \textbf{Dynamic CT and Digital Subtraction Myelography:} Fast, specialized myelogram techniques capture rapid dye flow to find tricky, fast-leaking spinal fluid tears and abnormal vein connections.
    \item \textbf{Metal-Reducing and Dual-Energy CT:} Advanced CT software reduces streaks around spine screws, rods, and cages, making postoperative nerve checks much sharper.
    \item \textbf{Spinal Fluid Lab Testing:} Fluid collected during the puncture can be checked for pressure, infection, inflammation, or abnormal cells in the same visit.
    \item \textbf{Nerve Studies (EMG and Nerve Conduction):} Electrical nerve tests check which pinched nerve seen on the myelogram is actually causing your pain or weakness, helping surgeons target the right spot.
\end{itemize}

\section{Comparison with Alternative Modalities}
\begin{itemize}
    \item \textbf{Compared With Spine MRI:} MRI is painless, uses no radiation or needle, and sees inside the spinal cord much better. CT myelography is essential when MRI is unsafe, blurred by surgical metal, or when movement views or spinal fluid leak mapping are needed.
    \item \textbf{Compared With Standard Spine CT:} A regular CT shows bone and hardware well, but cannot clearly outline the spinal cord or nerves. Adding contrast dye into the spinal fluid makes the nerves stand out sharply.
    \item \textbf{Compared With MR Myelography:} Special MRI sequences can highlight spinal fluid without a needle or radiation, but they cannot match dynamic CT myelography for pinpointing fast spinal fluid leaks.
\end{itemize}

\section{Follow-up Steps}
\begin{itemize}
    \item \textbf{Monitored Recovery and Fluids:} You will rest in a recovery area for 2--4 hours with your head propped up at 30--45 degrees while staff check for headache, nausea, bleeding, or new weakness. Drink plenty of fluids to help clear the dye and prevent headaches.
    \item \textbf{Resting and Restarting Medications:} Take it easy and avoid heavy lifting or strenuous exercise for 24--48 hours. Restart any paused blood thinners or other medications only when your doctor tells you it is safe.
    \item \textbf{Watching for Warning Signs:} Call your doctor or seek urgent care if you develop a severe headache that worsens when standing, fever, neck stiffness, new leg weakness, or trouble controlling your bladder or bowels. Mild post-puncture headaches usually improve with rest, fluids, caffeine, and pain relievers; stubborn ones can be treated with a simple procedure called an epidural blood patch.
    \item \textbf{Planning Next Treatment Steps:} Your spine specialist combines the myelogram, CT images, and any fluid lab results with your symptoms to plan targeted surgery, hardware adjustment, or sealing of a spinal fluid leak.
\end{itemize}

\section{What You Should Know About Myelography?}
\begin{itemize}
    \item \textbf{A Minor Spinal Procedure With Numbing Medicine:} Unlike a regular CT or MRI, myelography uses a thin needle to place dye into your spinal fluid. Local numbing medicine keeps discomfort to a brief sting and deep pressure, followed by a painless CT scan.
    \item \textbf{Radiation and Temporary Headache Risk:} Myelography uses X-ray radiation and specialized spinal dye. The most common side effect is a temporary headache when upright, which drinking extra water and resting afterward help prevent.
    \item \textbf{Plan for a Half-Day Visit:} The procedure and CT scan take about 45--60 minutes, followed by 2--4 hours of monitored recovery with your head slightly raised. Formal results usually reach your spine doctor within 1--2 business days.
    \item \textbf{Why Your Doctor Orders It:} Myelography gives surgeons crisp, reliable views of pinched nerves, bone spurs, surgical metal hardware, and hard-to-find spinal fluid leaks---especially when an MRI cannot be done or is blurred by metal.
    \item \textbf{What to Tell Your Care Team:} Always report blood thinners, past contrast dye allergies, seizure-risk medications, or possible pregnancy ahead of time. Arrange a ride home, and call right away afterward if you get a severe worsening headache, fever, neck stiffness, or new leg weakness.
\end{itemize}

\section{Key References}
\begin{itemize}
    \item American College of Radiology (ACR), American Society of Neuroradiology (ASNR), and Society for Pediatric Radiology (SPR). \textit{ACR--ASNR--SPR Practice Parameter for the Performance of Myelography and Cisternography}. Reston, VA: ACR; 2023.
    \item Patel IJ, Rahim S, Davidson JC, et al. Society of Interventional Radiology Consensus Guidelines for the Periprocedural Management of Thrombotic and Bleeding Risk in Patients Undergoing Percutaneous Image-Guided Interventions. \textit{Journal of Vascular and Interventional Radiology}. 2019;30(8):1168--1184.
    \item Farb RI, Nicholson PJ, Peng PW, et al. Spontaneous Intracranial Hypotension: A Systematic Imaging Approach for CSF Leak Localization and Management Based on MRI and Digital Subtraction Myelography. \textit{American Journal of Neuroradiology}. 2019;40(4):745--753.
    \item American College of Radiology (ACR). \textit{ACR Manual on Contrast Media}. Reston, VA: ACR; 2024.
\end{itemize}
